\documentclass[10pt,twocolumn,letterpaper]{article}

% ── Packages ──────────────────────────────────────────────────────────────────
\usepackage[
  top=0.75in, bottom=1in,
  left=0.625in, right=0.625in,
  columnsep=0.25in
]{geometry}
\usepackage{times}
\usepackage{amsmath,amssymb,amsfonts}
\usepackage{mathtools}
\usepackage{bm}
\usepackage{graphicx}
\usepackage{booktabs}
\usepackage{array}
\usepackage{caption}
\usepackage{xcolor}
\usepackage{flushend}
\usepackage{cite}
\usepackage{url}
\usepackage[hidelinks]{hyperref}
\usepackage{microtype}
\usepackage{titlesec}
\usepackage{abstract}
\usepackage{enumitem}

% ── IEEE-style section headings ───────────────────────────────────────────────
\titleformat{\section}
  {\normalfont\normalsize\bfseries\scshape\centering}
  {\Roman{section}.}{0.5em}{}
\titlespacing{\section}{0pt}{8pt plus 2pt minus 2pt}{4pt plus 1pt}

\titleformat{\subsection}
  {\normalfont\normalsize\itshape}
  {\Alph{subsection}.}{0.4em}{}
\titlespacing{\subsection}{0pt}{6pt plus 1pt minus 1pt}{3pt plus 1pt}

% ── Paragraph style ───────────────────────────────────────────────────────────
\setlength{\parindent}{1em}
\setlength{\parskip}{0pt}

% ── Caption style ─────────────────────────────────────────────────────────────
\captionsetup{font=small, labelfont=bf, labelsep=period}
\captionsetup[table]{position=above, labelfont={bf,sc}, justification=centering}

% ── Abstract override (full width, indented IEEE style) ───────────────────────
\renewenvironment{abstract}{%
  \twocolumn[{%
    % Title block
    \begin{center}
      \vspace*{6pt}
      {\fontsize{17}{20}\selectfont\bfseries
        Interpretable Multi-Scale Household Energy Forecasting via \\
        Holt--Winters Feature Augmentation, XGBoost, and Post-hoc Explanations
      }\par
      \vspace{12pt}
      {\normalsize
        Manaal Aamir \quad and \quad Arhum Khan \par
        \vspace{3pt}
        \textit{Department of Data Science, FAST University, Karachi, Pakistan} \par
        \vspace{2pt}
        \small manaal.aamir.2003@gmail.com \quad$|$\quad khan19@gmail.com
      }
      \vspace{14pt}
    \end{center}
    % Abstract box
    \begin{list}{}{%
      \setlength{\leftmargin}{0.45in}
      \setlength{\rightmargin}{0.45in}
      \setlength{\topsep}{0pt}
      \setlength{\itemsep}{0pt}
    }
    \item[]
    {\small\textbf{\textit{Abstract}---}\ignorespaces
}{%
    }
    \end{list}
    \vspace{4pt}
  }]
}

% ── Index terms environment ───────────────────────────────────────────────────
\newenvironment{IEEEkeywords}{%
  \begin{list}{}{%
    \setlength{\leftmargin}{0.45in}
    \setlength{\rightmargin}{0.45in}
    \setlength{\topsep}{0pt}
  }
  \item[]{\small\textbf{\textit{Index Terms}---}\ignorespaces}%
  \small
}{%
  \end{list}
  \vspace{4pt}
  \noindent\rule{\linewidth}{0.4pt}
  \vspace{6pt}
}

\pagestyle{empty}

% ─────────────────────────────────────────────────────────────────────────────
\begin{document}

\begin{abstract}
Most energy forecasting papers hand you a table with one RMSE per scale
and call it done. That number does not tell you when the model breaks
down or which features actually matter versus which just look important
on a SHAP plot. This paper builds a forecasting framework around that
gap. The approach combines Holt--Winters exponential smoothing with
XGBoost---an established pattern in the literature---but extends it with
leakage-aware feature engineering, temporal error decomposition by
hour-of-day, day-of-week, and month, and controlled feature intervention
studies through drop-one ablation and lag add-feature experiments. The
UCI Individual Household Electric Power Consumption dataset is used for
evaluation across hourly, daily, weekly, monthly, and quarterly
resolutions. Results show that aggregate RMSE hides non-uniform
degradation patterns only visible when error is sliced by time. Adding
autoregressive lag features materially improved performance at hourly
granularity and shifted SHAP attribution rankings---illustrating that
feature importance cannot be interpreted in isolation from what features
are present.
\end{abstract}

\begin{IEEEkeywords}
Energy forecasting, XGBoost, Holt--Winters, explainable AI, SHAP, LIME,
error analysis, feature ablation, time-series cross-validation.
\end{IEEEkeywords}

% ─────────────────────────────────────────────────────────────────────────────
\section{Introduction}

Residential electricity forecasting sits at the center of smart-grid
orchestration, demand response, and household energy optimization. Get
it wrong and operating costs rise, scheduling quality drops, and users
stop trusting the automated systems built on top of the forecasts. Get
it right on the aggregate metric but miss the temporal failure patterns,
and the problem is roughly the same---you just do not know it yet.

Gradient-boosted tree models, XGBoost in particular, consistently
achieve strong predictive performance on tabular time-series problems.
The trade-off is interpretability. A model that a domain user or system
operator cannot interrogate is difficult to trust, regardless of its
error metrics. Post-hoc explanation methods---SHAP, LIME, permutation
importance, partial dependence plots---partially address this, but they
report attributions given a fixed feature set. They cannot tell you
whether a highly attributed feature is load-bearing or whether removing
it would barely move the needle.

The pipeline built in this work starts from an existing
approach---Holt--Winters decomposition for feature augmentation fed into
XGBoost---and extends it in three concrete directions the base
formulation does not cover. First, leakage in the HW component features
is fixed by lagging them one step, so that predicting $y$ at $t{+}1$
uses only values known at $t$. Second, the RMSE is broken into time
slices by hour-of-day, day-of-week, and month, exposing failure windows
a single number cannot reveal. Third, feature interventions---drop-one
ablation and lag add-feature experiments---measure accuracy changes
directly rather than relying on attribution alone.

\subsection{Contributions}

\begin{itemize}[leftmargin=1.5em, itemsep=2pt, topsep=3pt]
  \item A reproducible HW+XGB+XAI pipeline for household energy
    forecasting across hourly, daily, weekly, monthly, and quarterly
    resolutions, runnable end-to-end from a single CLI command.
  \item Leakage-aware temporal modeling via chronological splits,
    \texttt{TimeSeriesSplit} cross-validation, and one-step-lagged
    Holt--Winters component features.
  \item Error decomposition by hour-of-day (24 buckets), day-of-week
    (7 buckets), and month (12 buckets), exposing non-uniform
    degradation patterns hidden by aggregate RMSE.
  \item Feature intervention analysis via drop-one ablation and
    target-lag add-feature experiments, quantifying accuracy deltas and
    corresponding SHAP attribution shifts.
\end{itemize}

% ─────────────────────────────────────────────────────────────────────────────
\section{Literature Review}

\subsection{Household Energy Forecasting}

Household electricity forecasting methods split into three broad groups:
classical statistical models such as exponential smoothing and ARIMA
variants, machine learning approaches built around tree ensembles, and
deep learning architectures~\cite{hyndman2021}. Statistical models
handle seasonality well but underfit non-linear demand patterns. Deep
learning models capture complex dependencies but require substantial
data and are difficult to explain to system operators. Tree ensembles
sit in the middle: strong predictive accuracy, moderate training cost,
and amenable to post-hoc explanation~\cite{chen2016xgboost}.

Hybrid approaches combining classical decomposition with machine
learning are an active area. Holt--Winters exponential smoothing is a
common decomposition choice because it handles both trend and
seasonality with few parameters and is computationally cheap. Most
published HW+ML pipelines either forecast each component separately
and combine outputs, or feed component estimates as features to a
downstream learner. This work takes the latter approach.

\subsection{Interpretable Machine Learning for Energy Systems}

Post-hoc interpretability methods divide into global and local
explanations. SHAP provides additive feature attributions grounded in
cooperative game theory, with efficient TreeSHAP implementations for
gradient-boosted models~\cite{lundberg2017shap}. LIME fits a local
linear surrogate around a specific prediction~\cite{ribeiro2016lime}.
Permutation Feature Importance (PFI) shuffles one feature at a time
and measures the accuracy drop. Partial Dependence Plots (PDP) show
the marginal effect of a feature averaged over the distribution of
other inputs~\cite{goldstein2015peeking}.

Using multiple methods together is more informative than any single
one. SHAP and PFI do not always agree, and that disagreement often
indicates feature interactions or instability under perturbation---both
worth reporting rather than explaining away.

\subsection{Research Gap}

The standard evaluation practice in energy forecasting is to report a
single aggregate metric per scale and present SHAP or PFI plots as
evidence of interpretability. This has two problems. First, aggregate
RMSE obscures temporal failure concentration. Second, feature
attribution without feature intervention does not establish causal
importance---a feature can rank highly on SHAP while being almost
entirely redundant. This work addresses both gaps directly.

% ─────────────────────────────────────────────────────────────────────────────
\section{Dataset and Preprocessing}

\subsection{Dataset}

All experiments use the UCI Individual Household Electric Power
Consumption dataset~\cite{uci2012}, recorded at one-minute resolution
over roughly four years from a single household in Sceaux, France. The
dataset includes global active power $P_{\text{global}}$, reactive
power $Q_{\text{global}}$, voltage $V$, global intensity $I$, and
three sub-metering channels $S_1, S_2, S_3$ corresponding to the
kitchen, laundry room, and water heater or air conditioning loads.

\subsection{Cleaning and Feature Construction}

Let $\mathcal{T} = \{t_1, t_2, \ldots, t_N\}$ denote the full
timestamp index. Missing values are filled using time-based linear
interpolation:
\begin{equation}
  x(t) = x(t_a) + \frac{t - t_a}{t_b - t_a}\bigl(x(t_b) - x(t_a)\bigr),
  \label{eq:interp}
\end{equation}
where $t_a$ and $t_b$ are the nearest non-missing timestamps bracketing
$t$, followed by forward and backward fill for residual gaps.

Cyclical time encodings are added to avoid boundary discontinuities
with raw integer features. For hour-of-day $h \in \{0,\ldots,23\}$
and day-of-week $d \in \{0,\ldots,6\}$:
\begin{align}
  \text{hour\_sin} &= \sin\!\left(\tfrac{2\pi h}{24}\right), &
  \text{hour\_cos} &= \cos\!\left(\tfrac{2\pi h}{24}\right), \label{eq:hourcyc}\\[2pt]
  \text{dow\_sin}  &= \sin\!\left(\tfrac{2\pi d}{7}\right),  &
  \text{dow\_cos}  &= \cos\!\left(\tfrac{2\pi d}{7}\right).  \label{eq:dowcyc}
\end{align}
These encodings give the model a smooth periodic representation of
time: Euclidean distance between any two time points is proportional
to their circular distance on the clock or week cycle.

\subsection{Multi-Scale Aggregation}

The one-minute series is resampled to five temporal granularities
$\mathcal{G} = \{\text{hourly, daily, weekly, monthly, quarterly}\}$.
For a time window $\mathcal{W}$ of width $g \in \mathcal{G}$:
\begin{equation}
  \tilde{x}^{(g)}_{\mathcal{W}} =
  \begin{cases}
    \dfrac{1}{|\mathcal{W}|}\displaystyle\sum_{t \in \mathcal{W}} x_t
      & \text{power-like variables,}\\[6pt]
    \displaystyle\sum_{t \in \mathcal{W}} x_t
      & \text{sub-metering channels.}
  \end{cases}
  \label{eq:resample}
\end{equation}
Summing the energy channels preserves total energy consumed per period.
A stationarity check via the Augmented Dickey--Fuller (ADF) test is
run on the training target per scale:
\begin{equation}
  \tau_{\text{ADF}} = \frac{\hat{\delta}}{\operatorname{SE}(\hat{\delta})},
  \label{eq:adf}
\end{equation}
where $\hat{\delta}$ is the OLS estimate of the unit-root coefficient
in $\Delta y_t = \alpha + \beta t + \delta y_{t-1}
+ \sum_{k=1}^{p}\gamma_k\Delta y_{t-k} + \varepsilon_t$.
Results are written to JSON as a diagnostic only.

% ─────────────────────────────────────────────────────────────────────────────
\section{Methodology}

\subsection{Holt--Winters Feature Augmentation}

The additive Holt--Winters model decomposes the training target $y_t$:
\begin{equation}
  y_t = \ell_t + b_t + s_t + \varepsilon_t,
  \label{eq:hw_decomp}
\end{equation}
where $\ell_t$ is the level, $b_t$ the trend, $s_t$ the seasonal
component with period $m$, and $\varepsilon_t$ the residual. The
smoothing recursions are:
\begin{align}
  \ell_t &= \alpha(y_t - s_{t-m}) + (1{-}\alpha)(\ell_{t-1} + b_{t-1}),\label{eq:hw_level}\\
  b_t    &= \beta^*(\ell_t - \ell_{t-1}) + (1{-}\beta^*)b_{t-1},\label{eq:hw_trend}\\
  s_t    &= \gamma(y_t - \ell_{t-1} - b_{t-1}) + (1{-}\gamma)s_{t-m},\label{eq:hw_season}
\end{align}
with $\alpha,\beta^*,\gamma \in [0,1]$ estimated by minimising
sum-of-squared one-step forecast errors on training data:
\begin{equation}
  \min_{\alpha,\beta^*,\gamma}\;
  \sum_{t=1}^{T_{\text{train}}} \bigl(y_t - \hat{y}_{t|t-1}\bigr)^2.
  \label{eq:hw_optim}
\end{equation}

The standard HW+ML formulation suffers from leakage: component
estimates at $t$ depend on $y_t$ yet are used to predict $y_t$.
The fix applied here lags all components by one step:
\begin{equation}
  \bm{\phi}_t =
  \bigl[\hat{\ell}_{t-1},\;\hat{b}_{t-1},\;\hat{s}_{t-1},\;\hat{\varepsilon}_{t-1}\bigr]^\top,
  \label{eq:hw_lag}
\end{equation}
so that predicting $y_{t+1}$ uses only component values observable
at~$t$.

\subsection{Supervised Learning Formulation}

The full feature vector at time $t$ is:
\begin{equation}
  \mathbf{x}_t =
  \bigl[\mathbf{z}_t^\top,\;\mathbf{c}_t^\top,\;\bm{\phi}_t^\top,\;
  (\mathbf{l}_t^\top)\bigr]^\top,
  \label{eq:feature_vec}
\end{equation}
where $\mathbf{z}_t$ are original covariates,
$\mathbf{c}_t = [\text{hour\_sin},\text{hour\_cos},\text{dow\_sin},\text{dow\_cos}]^\top$,
$\bm{\phi}_t$ are the lagged HW components from~\eqref{eq:hw_lag},
and $\mathbf{l}_t$ is an optional autoregressive lag block:
\begin{equation}
  \mathbf{l}_t = \bigl[y_{t-1},\;y_{t-24},\;y_{t-168}\bigr]^\top.
  \label{eq:lag_features}
\end{equation}
The prediction target is $y_{t+1}$:
\begin{equation}
  \hat{y}_{t+1} = f(\mathbf{x}_t).
  \label{eq:supervised}
\end{equation}

\subsection{XGBoost Regression and Tuning}

XGBoost learns an additive ensemble of $K$ regression trees:
\begin{equation}
  \hat{y}_{t+1} = \sum_{k=1}^{K} f_k(\mathbf{x}_t),\quad f_k \in \mathcal{F}.
  \label{eq:xgb_ensemble}
\end{equation}
The objective at iteration $k$ is:
\begin{equation}
  \mathcal{L}^{(k)} =
  \sum_{t} \ell\!\left(y_{t+1},\,\hat{y}_{t+1}^{(k-1)} + f_k(\mathbf{x}_t)\right)
  + \Omega(f_k),
  \label{eq:xgb_obj}
\end{equation}
with regularisation:
\begin{equation}
  \Omega(f_k) = \gamma T + \tfrac{1}{2}\lambda\|\mathbf{w}\|^2,
  \label{eq:xgb_reg}
\end{equation}
where $T$ is the number of leaves, $\mathbf{w}$ the leaf weight
vector, $\gamma$ the minimum split gain, and $\lambda$ the $L_2$
penalty. Hyperparameter tuning uses \texttt{TimeSeriesSplit}
($n{=}5$) and \texttt{GridSearchCV}. Feature standardisation:
\begin{equation}
  \tilde{x}_{tj} =
  \frac{x_{tj} - \mu_j^{\text{train}}}{\sigma_j^{\text{train}}},
  \label{eq:standardise}
\end{equation}
with $\mu_j^{\text{train}}, \sigma_j^{\text{train}}$ computed on
training data only.

% ─────────────────────────────────────────────────────────────────────────────
\section{Explainability and Diagnostics}

\subsection{Post-hoc Explanation Methods}

Four complementary methods are applied to each trained model.

\textbf{SHAP (TreeSHAP)} computes Shapley values:
\begin{equation}
  \phi_j =
  \sum_{S \subseteq \mathcal{F}\setminus\{j\}}
  \frac{|S|!\,(|\mathcal{F}|-|S|-1)!}{|\mathcal{F}|!}
  \bigl[v(S\cup\{j\}) - v(S)\bigr],
  \label{eq:shap}
\end{equation}
where $v(S)$ is the model output conditioned on features in
$S$~\cite{lundberg2017shap}. TreeSHAP runs in
$\mathcal{O}(TLD^2)$ time.

\textbf{LIME} fits a local linear surrogate $g$ around instance
$\mathbf{x}_0$:
\begin{equation}
  g^* = \operatorname*{arg\,min}_{g\in\mathcal{G}}
  \sum_{z}\pi_{\mathbf{x}_0}(z)\,\bigl[f(z)-g(z)\bigr]^2 + \Omega(g),
  \label{eq:lime}
\end{equation}
with proximity kernel
$\pi_{\mathbf{x}_0}(z) = \exp\!\left(-D(z,\mathbf{x}_0)^2/\sigma^2\right)$~\cite{ribeiro2016lime}.

\textbf{Permutation Feature Importance (PFI)}:
\begin{equation}
  \text{PFI}_j =
  \text{RMSE}(\mathbf{X}_{\pi_j}) - \text{RMSE}(\mathbf{X}),
  \label{eq:pfi}
\end{equation}
where $\mathbf{X}_{\pi_j}$ has column $j$ randomly permuted.

\textbf{Partial Dependence (PDP)} estimates the marginal effect of
feature $j$:
\begin{equation}
  \hat{f}_j(x_j) =
  \frac{1}{n}\sum_{i=1}^{n} f\!\left(x_j,\,\mathbf{x}_i^{(-j)}\right).
  \label{eq:pdp}
\end{equation}

\subsection{Error Decomposition by Time Slice}

For each temporal bucket $b$ in $\{\text{hour-of-day, day-of-week,
month}\}$, the slice RMSE is:
\begin{equation}
  \text{RMSE}_{b} =
  \sqrt{\frac{1}{|\mathcal{I}_b|}\sum_{t \in \mathcal{I}_b}
  \bigl(y_{t+1} - \hat{y}_{t+1}\bigr)^2},
  \label{eq:slice_rmse}
\end{equation}
where $\mathcal{I}_b = \{t \in \mathcal{T}_{\text{test}} :
\text{bucket}(t) = b\}$. The 24 hourly, 7 daily, and 12 monthly
slices together produce a map of where the model degrades---invisible
in the scalar aggregate.

\subsection{Feature Impact Studies}

\textbf{Drop-one ablation.} For each feature $j$ in the top-$N$ by
XGBoost gain, retrain on $\mathcal{F}\setminus\{j\}$ and compute:
\begin{align}
  \Delta\text{RMSE}_j &=
    \text{RMSE}(f_{\mathcal{F}\setminus\{j\}}) - \text{RMSE}(f_{\mathcal{F}}),
    \label{eq:ablation_rmse}\\
  \Delta R^2_j &=
    R^2(f_{\mathcal{F}}) - R^2\!\left(f_{\mathcal{F}\setminus\{j\}}\right).
    \label{eq:ablation_r2}
\end{align}
Positive $\Delta\text{RMSE}_j$ indicates the feature was genuinely
load-bearing; zero or negative values suggest redundancy.

\textbf{Add-feature experiment.} Let
$\mathcal{F}^+ = \mathcal{F} \cup \mathbf{l}_t$ be the augmented
feature set. Record:
\begin{equation}
  \Delta\text{RMSE}^+ =
  \text{RMSE}(f_{\mathcal{F}}) - \text{RMSE}(f_{\mathcal{F}^+}),
  \label{eq:add_feat}
\end{equation}
alongside SHAP ranking shifts between $f_{\mathcal{F}}$ and
$f_{\mathcal{F}^+}$.

% ─────────────────────────────────────────────────────────────────────────────
\section{Experimental Setup}

\subsection{Evaluation Metrics}

Four metrics are reported for test set $\{(y_i,\hat{y}_i)\}_{i=1}^n$:
\begin{align}
  \text{MAE}  &= \tfrac{1}{n}{\textstyle\sum}|y_i - \hat{y}_i|,
    \label{eq:mae}\\
  \text{MSE}  &= \tfrac{1}{n}{\textstyle\sum}(y_i - \hat{y}_i)^2,
    \label{eq:mse}\\
  \text{RMSE} &= \sqrt{\tfrac{1}{n}{\textstyle\sum}(y_i-\hat{y}_i)^2},
    \label{eq:rmse}\\
  R^2         &= 1 - \frac{{\textstyle\sum}(y_i-\hat{y}_i)^2}
                          {{\textstyle\sum}(y_i-\bar{y})^2}.
    \label{eq:r2}
\end{align}
RMSE is the primary comparison metric. $R^2$ provides a
scale-independent reference. MAE is reported as a robustness check.

\subsection{Train/Test Protocol}

An 80/20 chronological split is used throughout; no shuffling occurs
at any stage. Hyperparameter selection uses five-fold
\texttt{TimeSeriesSplit} cross-validation applied exclusively to the
training portion. The test set is touched once only, after
hyperparameters are selected, to compute final metrics.

% ─────────────────────────────────────────────────────────────────────────────
\section{Results}

\subsection{Aggregate Performance Across Scales}

Table~\ref{tab:metrics} reports MAE, RMSE, and $R^2$ for each
temporal scale. Performance generally improves at coarser
scales---aggregation smooths short-term variability that is difficult
to predict from available covariates alone. The hourly scale is the
hardest, with the most volatile target and the strongest dependency
on recent history.

The negative $R^2$ at daily scale ($-0.2280$) is a clear signal: the
model underperforms a naive mean predictor at that granularity. This
likely reflects the smoothing effect of daily aggregation reducing
the variance the model can exploit, combined with insufficient
feature information to capture day-to-day fluctuations. Monthly
($R^2 = 0.7580$) and quarterly ($R^2 = 0.8127$) scales show
substantially stronger fit, where seasonal and trend components
dominate and are well-captured by the HW augmentation.

\begin{table}[t]
  \caption{Forecasting Performance by Temporal Scale}
  \label{tab:metrics}
  \centering
  \setlength{\tabcolsep}{6pt}
  \begin{tabular}{lrrr}
    \toprule
    \textbf{Scale} & \textbf{MAE} & \textbf{RMSE} & $\bm{R^2}$ \\
    \midrule
    Hourly    & 0.4002 & 0.5258 & \phantom{$-$}0.4787 \\
    Daily     & 0.2842 & 0.3457 & $-$0.2280           \\
    Weekly    & 0.1432 & 0.1977 & \phantom{$-$}0.3356 \\
    Monthly   & 0.0902 & 0.1037 & \phantom{$-$}0.7580 \\
    Quarterly & 0.0627 & 0.0721 & \phantom{$-$}0.8127 \\
    \bottomrule
  \end{tabular}
  \par\smallskip
  {\footnotesize Source: \texttt{outputs/tables/metrics\_latest.csv}.
   Negative $R^2$ at daily scale indicates the model underperforms
   a naive mean predictor at that granularity.}
\end{table}

\subsection{Temporal Error Decomposition}

Slicing test RMSE by hour-of-day, day-of-week, and month using
\eqref{eq:slice_rmse} confirms that aggregate metrics hide
operationally important variation. Evening hours, when household
demand peaks and usage patterns are more variable, consistently show
elevated RMSE relative to overnight windows. Day-of-week slices
reveal higher error on weekdays at certain scales, reflecting
work-schedule-dependent consumption patterns the model captures
imperfectly. Monthly slices expose seasonal variation: winter months
show higher absolute error where heating or cooling loads dominate.
These patterns indicate where confidence intervals should be widened
and where additional features---weather, occupancy proxies---would
provide the most benefit.

\subsection{Feature Ablation and Lag Add-Feature Results}

Drop-one ablation, computed via
\eqref{eq:ablation_rmse}--\eqref{eq:ablation_r2}, shows that not all
top-SHAP features are equally load-bearing. Several features that
ranked highly by gain showed only modest $\Delta\text{RMSE}$ when
removed, suggesting shared predictive contribution with correlated
features. A smaller subset showed clear positive $\Delta\text{RMSE}$
on removal, confirming genuine predictive necessity.

Adding target lag features via \eqref{eq:lag_features} at hourly
resolution improved RMSE materially, with $\Delta\text{RMSE}^+$ from
\eqref{eq:add_feat} driven primarily by the 1-step and 24-step lags.
The SHAP ranking shifted after lags were added: HW-derived features
dropped in relative attribution because the autoregressive lags
absorbed much of the short-term predictive load. This demonstrates
why feature importance rankings must not be interpreted independently
of what features are present.

\subsection{Explainability Outcomes}

SHAP summary plots show global feature ranking and directionality.
The HW level component $\hat{\ell}_{t-1}$ consistently ranks among
the top features---level tracks the current baseline of consumption,
and the next-step value is closely tied to it. Cyclical time features
(hour\_sin, hour\_cos) rank highly at hourly scale and less so at
coarser scales where within-day variation is averaged out.

LIME explanations for selected test instances show that local feature
rankings can differ substantially from global SHAP rankings,
particularly for unusual consumption periods. PFI rankings computed
via \eqref{eq:pfi} do not always match SHAP rankings---a signal of
features whose importance depends on interaction effects rather than
independent marginal contribution. PDP plots via \eqref{eq:pdp}
confirm expected monotonic relationships for level and trend
components and nonlinear response curves for cyclical time features.

% ─────────────────────────────────────────────────────────────────────────────
\section{Discussion}

Three practical points stand out.

First, aggregate metrics are not enough for deployment decisions. A
model that looks fine on average can be substantially degraded during
the hours and seasons that matter most. Temporal error decomposition
\eqref{eq:slice_rmse} is cheap to compute and adds real operational
value---it should be a default part of any evaluation pipeline, not
an optional diagnostic.

Second, feature attribution is not the same as feature necessity.
High Shapley values \eqref{eq:shap} are often shared across
correlated features. Ablation studies
\eqref{eq:ablation_rmse}--\eqref{eq:ablation_r2} are the more direct
test, and their results do not always match attribution rankings.
Both should be reported when interpretability matters for deployment.

Third, the HW+XGB combination offers a reasonable middle ground
between fully statistical and fully black-box approaches. The
Holt--Winters recursions \eqref{eq:hw_level}--\eqref{eq:hw_season}
are transparent and produce interpretable components. XGBoost handles
the non-linear residuals via \eqref{eq:xgb_ensemble}. The combination
is not the highest-accuracy possible approach, but it is
substantially more interpretable and easier to audit.

\subsection{Threats to Validity}

The main limitations are single-dataset evaluation, limited
hyperparameter search breadth, and absence of uncertainty
quantification. The dataset covers a single household in France;
generalisation to different household types, climates, or grid
contexts is an open question. The daily-scale negative $R^2$ warrants
further investigation: it may reflect aggregation collapsing the
predictive signal present at finer granularities.

% ─────────────────────────────────────────────────────────────────────────────
\section{Conclusion and Future Work}

This paper built and evaluated an interpretable multi-scale energy
forecasting pipeline combining Holt--Winters feature augmentation
\eqref{eq:hw_decomp}--\eqref{eq:hw_lag}, XGBoost regression
\eqref{eq:xgb_ensemble}, and four post-hoc explanation methods. The
core prediction approach existed in prior work. What was added---
leakage-aware component lagging, temporal error decomposition, drop-one
ablation, and lag add-feature experiments---addresses the evaluation
and interpretability gaps most published HW+ML work leaves open.

The main result is not a new state-of-the-art accuracy number. It is
a more honest picture of where a well-built baseline model succeeds and
where it struggles, and a cleaner methodology for connecting feature
importance claims to actual predictive necessity. The daily-scale
result ($R^2 = -0.228$) is a useful reminder that aggregation does
not always help: smoothing can collapse the very signal a model needs.

Future work will focus on: (i)~cross-dataset generalisation to
evaluate robustness across household types and geographies;
(ii)~uncertainty quantification via conformal prediction or quantile
regression; (iii)~concept-drift detection and adaptive retraining;
and (iv)~richer contextual features---weather data, occupancy
proxies---most likely to reduce error in the peak-period windows
identified by the error decomposition analysis.

% ─────────────────────────────────────────────────────────────────────────────
\section*{Reproducibility}

All experiments run end-to-end from a single CLI entry point. Outputs
are organised under \texttt{outputs/} as follows:
\texttt{metrics\_latest.csv} and \texttt{metrics\_history.csv} for
performance; \texttt{error\_slice\_*} tables and plots for temporal
decomposition; \texttt{ablation\_drop1\_*} for ablation;
\texttt{pfi\_*} for permutation importance; \texttt{shap\_summary\_*}
and \texttt{pdp\_*} for global explanations; and \texttt{lime\_*}
HTML for per-instance local explanations.

% ─────────────────────────────────────────────────────────────────────────────
\begin{thebibliography}{00}

\bibitem{uci2012}
UCI Machine Learning Repository, ``Individual Household Electric Power
Consumption Dataset,'' 2012. [Online]. Available:
\url{https://archive.ics.uci.edu/ml/datasets/individual+household+electric+power+consumption}

\bibitem{chen2016xgboost}
T.~Chen and C.~Guestrin, ``XGBoost: A Scalable Tree Boosting System,''
in \textit{Proc.\ 22nd ACM SIGKDD Int.\ Conf.\ Knowledge Discovery and
Data Mining}, San Francisco, CA, 2016, pp.~785--794.

\bibitem{lundberg2017shap}
S.~M.~Lundberg and S.-I.~Lee, ``A Unified Approach to Interpreting
Model Predictions,'' in \textit{Advances in Neural Information
Processing Systems (NeurIPS)}, vol.~30, 2017.

\bibitem{ribeiro2016lime}
M.~T.~Ribeiro, S.~Singh, and C.~Guestrin, ```Why Should I Trust
You?': Explaining the Predictions of Any Classifier,'' in
\textit{Proc.\ 22nd ACM SIGKDD Int.\ Conf.\ Knowledge Discovery and
Data Mining}, 2016, pp.~1135--1144.

\bibitem{hyndman2021}
R.~J.~Hyndman and G.~Athanasopoulos, \textit{Forecasting: Principles
and Practice}, 3rd~ed. Melbourne, Australia: OTexts, 2021. [Online].
Available: \url{https://otexts.com/fpp3}

\bibitem{holt2004}
C.~C.~Holt, ``Forecasting Seasonals and Trends by Exponentially
Weighted Moving Averages,'' \textit{Int.\ J.\ Forecasting}, vol.~20,
no.~1, pp.~5--10, 2004.

\bibitem{winters1960}
P.~R.~Winters, ``Forecasting Sales by Exponentially Weighted Moving
Averages,'' \textit{Management Science}, vol.~6, no.~3,
pp.~324--342, 1960.

\bibitem{goldstein2015peeking}
A.~Goldstein, A.~Kapelner, J.~Bleich, and E.~Pitkin, ``Peeking Inside
the Black Box: Visualizing Statistical Learning with Plots of
Individual Conditional Expectation,'' \textit{J.\ Computational and
Graphical Statistics}, vol.~24, no.~1, pp.~44--65, 2015.

\end{thebibliography}

\end{document}